% Chapter 17: FX and multi-asset options.
\providecommand{\Corr}{\operatorname{Corr}}
\providecommand{\Black}{\operatorname{Black}}

A currency option is the simplest option on two assets: the foreign money market and the
domestic one. Its price depends on which currency we count in, and the answer to ``what is it
worth?'' must not. This invariance is a change of \emph{num\'eraire}, the unit in which prices
are counted, and it is the one tool that runs through the chapter. It prices currency options
from either side of the market (\cref{sec:ma-fx,sec:ma-numeraire}), turns the option to
exchange one asset for another into a Black--Scholes call (\cref{sec:ma-exchange}), and
explains the drift adjustment of quanto contracts (\cref{sec:ma-quanto}). Payoffs that do not
reduce to one dimension, such as baskets, are priced by simulating correlated assets, with the
geometric basket as a control variate, as for the Asian options of
\cref{ch:15-asian} (\cref{sec:ma-basket}).

We write calls on lognormal variables with \emph{Black's function}: for $F>0$, $v>0$ and $K>0$,
\begin{equation}\label{eq:ma-black}
  \Black(F,K,v)=F\,N(d_+)-K\,N(d_-),\qquad d_\pm=\frac{\log(F/K)\pm\frac12v}{\sqrt v},
\end{equation}
and $\Black(F,K,0)=(F-K)^+$, $\Black(F,K,v)=F-K$ for $K\le0$. If $L$ is lognormal with mean
$F$ and $\Var\log L=v$, then $\E(L-K)^+=\Black(F,K,v)$: this is the computation behind
\cref{thm:bsf-formula}, and the Black--Scholes call is
$\mathrm{BS}(S,K,T;r,y,\sigma)=e^{-rT}\Black(Se^{(r-y)T},K,\sigma^2T)$. Black's function is
positively homogeneous of degree one in $(F,K)$.

\section{Currency options}\label{sec:ma-fx}

A currency market has a domestic and a foreign money market, with constant rates $r_d$ and
$r_f$, and an \emph{exchange rate} $X_t$, the domestic price of one unit of foreign currency.
With the dollar domestic and the pound foreign, $X$ is the USD price of one GBP, quoted as
GBP/USD (``cable''). A pound deposited at the rate $r_f$ is worth $X_te^{r_ft}$ dollars at time
$t$, so for a dollar investor the pound is an asset that pays a continuous dividend at the rate
$r_f$. Under a risk-neutral measure $\Q$ for that investor (\cref{prop:sm-risk-neutral}),
$e^{-r_dt}X_te^{r_ft}$ is a martingale. The model of \textcite{garman-kohlhagen} is the
standard model of \cref{def:sm-model} for this asset:
\begin{equation}\label{eq:ma-gk}
  \frac{\dd X_t}{X_t}=(r_d-r_f)\,\dd t+\sigma\,\dd W_t\quad\text{under }\Q,\qquad\sigma>0,
\end{equation}
that is, $S=X$, $r=r_d$ and $y=r_f$: the foreign rate is the yield.

\begin{proposition}[Garman--Kohlhagen]\label{prop:ma-garman-kohlhagen}
In the model \eqref{eq:ma-gk}, let $K>0$, $T>0$, $F=X_0e^{(r_d-r_f)T}$ be the forward
exchange rate and $d_\pm=\big(\log(F/K)\pm\frac12\sigma^2T\big)/(\sigma\sqrt T)$.
\begin{enumerate}[label=(\roman*)]
  \item A call on one unit of foreign currency with strike $K$, paying $(X_T-K)^+$ units of
    domestic currency, and the put, paying $(K-X_T)^+$, cost in domestic currency
    \begin{align*}
      C&=X_0e^{-r_fT}N(d_+)-Ke^{-r_dT}N(d_-)=e^{-r_dT}\Black(F,K,\sigma^2T),\\
      P&=Ke^{-r_dT}N(-d_-)-X_0e^{-r_fT}N(-d_+),
    \end{align*}
    and $C/X_0$, $P/X_0$ in foreign currency. Parity reads $C-P=e^{-r_dT}(F-K)$; at the
    money forward, $K=F$, $C=P=e^{-r_dT}F\big(2N(\frac12\sigma\sqrt T)-1\big)$. If $r_d=r_f$, then
    $F=X_0$ and at the money forward is at the money spot.
  \item The replicating portfolio of the call holds, at time $0$,
    $\Delta_f=\partial C/\partial X_0=e^{-r_fT}N(d_+)$ units of foreign currency and
    $\Delta_d=C-X_0\,\partial C/\partial X_0=-Ke^{-r_dT}N(d_-)$ units of domestic currency, both
    in the money markets; $C=X_0\Delta_f+\Delta_d$.
\end{enumerate}
\end{proposition}

\begin{proof}
(i) is \cref{thm:bsf-formula} with $S=X$, $r=r_d$, $y=r_f$; the ATMF value uses $d_\pm=\pm\frac12\sigma\sqrt T$.
(ii) The delta is that of \cref{prop:bsf-greeks}; it is a number of units of the asset, here
foreign currency, held with its ``dividends'' (the foreign interest) reinvested, as in the
replication of \cref{prop:sm-risk-neutral}(ii). The rest of the portfolio is cash,
$C-X_0\Delta_f=-Ke^{-r_dT}N(d_-)$. The price is homogeneous of degree one in $(X_0,K)$, so
$\Delta_d=K\,\partial C/\partial K$ by Euler's theorem (\cref{prop:sr-euler}).
\end{proof}

\begin{example}[The course's cable call]\label{exa:ma-cable}
With $X_0=1.2000$, $r_d=r_{\mathrm{USD}}=3\%$, $r_f=r_{\mathrm{GBP}}=2.5\%$, $\sigma=11\%$ and $T=1$, the forward
is $F=1.2\,e^{0.005}=1.206015$, and the ATMF call costs $C=0.0513343$ USD, that is
$0.0427786$ GBP, or $4.2565\%$ of the strike. Its hedge holds $\Delta_f=0.509044$ GBP and
$\Delta_d=-0.559519$ USD, and indeed $1.2\times0.509044-0.559519=0.051334$. These are the
course's $0.05133$, $0.04278$, $0.5090$ and $-0.5595$.
\end{example}

Currency markets quote differently from equity markets. Options are struck relative to the
forward (ATMF rather than at the spot), premiums are quoted in either currency, and implied
volatilities are quoted against a delta rather than a strike. The course's GBP/USD and
EUR/USD smiles (CME settlement prices, 22 July 2022), plotted against the normalised strike
$\log(K/F)/(\sigma_{\mathrm{atm}}\sqrt T)$, fall from low strikes to a minimum above the forward and
then rise again: a skew, as for equity indices (\cref{rem:sk-smile-data}), with a smile on the upside.
Neither currency is special: by \cref{prop:ma-numeraire}(v) below, the USD/GBP smile is the mirror
image of the GBP/USD smile, $\sigma_{\mathrm{USD/GBP}}(m)=\sigma_{\mathrm{GBP/USD}}(-m)$ in the
log-moneyness $m=\log(K/F)$.

\section{Changing the num\'eraire}\label{sec:ma-numeraire}

In this section $W$ is a $d$-dimensional Brownian motion under $\Q$ on $[0,T]$, $(\F_t)$ is its
augmented natural filtration, and $B_t=e^{rt}$ is the money market. A \emph{price process} is
an adapted process $V$ such that $V/B$ is a $\Q$-martingale on $[0,T]$: the value of a traded
asset that pays no dividends, of a self-financing portfolio, or of a claim. A stock that pays
the dividend yield $y$ is not one, but $e^{yt}S_t$, one share at $0$ with its dividends
reinvested, is. Counting prices in units of a positive price process $N$ instead of money is a
\emph{change of num\'eraire}; \cref{rem:sk-calendar} already counted prices in prepaid forwards.

\begin{proposition}[Change of num\'eraire]\label{prop:ma-numeraire}
Let $N$ be a price process with $N_t>0$ for all $t\le T$, almost surely.
\begin{enumerate}[label=(\roman*)]
  \item $Z_T=\dfrac{N_TB_0}{N_0B_T}$ is positive with $\EQ Z_T=1$, so
    $\dd\Q^N=Z_T\,\dd\Q$ defines a probability measure on $\F_T$ equivalent to $\Q$, and
    $\EQ[Z_T\mid\F_t]=N_tB_0/(N_0B_t)$.
  \item (\emph{Num\'eraire invariance.}) An adapted process $V$ is a price process if and only
    if $V/N$ is a $\Q^N$-martingale. For every $\F_T$-measurable $H$ with
    $\EQ[|H|/B_T]<\infty$, equivalently $\E^{\Q^N}[|H|/N_T]<\infty$,
    \[
      B_t\,\EQ\Big[\frac{H}{B_T}\,\Big|\,\F_t\Big]=N_t\,\E^{\Q^N}\Big[\frac{H}{N_T}\,\Big|\,\F_t\Big],
      \qquad0\le t\le T .
    \]
  \item (\emph{Drifts.}) Suppose $N_t/B_t=N_0\exp(\theta\cdot W_t-\frac12|\theta|^2t)$ for a
    constant $\theta\in\R^d$. Then $W^N_t=W_t-\theta t$ is a $\Q^N$-Brownian motion. If
    $\dd S_t/S_t=\mu_t\,\dd t+\sigma_S\cdot\dd W_t$ with $\sigma_S\in\R^d$, then
    $\dd S_t/S_t=(\mu_t+\sigma_S\cdot\theta)\,\dd t+\sigma_S\cdot\dd W^N_t$: the drift rises by
    $\sigma_S\cdot\theta$, the instantaneous covariance of $\log S$ and $\log N$, and the
    volatility does not change.
  \item (\emph{The share measure.}) If $\dd S_t/S_t=(r-y)\,\dd t+\sigma_S\cdot\dd W_t$ with
    constants $y$ and $\sigma_S$, then $N_t=e^{yt}S_t$ satisfies (iii) with $\theta=\sigma_S$,
    \[
      \frac{\dd\Q^S}{\dd\Q}=\frac{S_T}{\EQ[S_T]},
    \]
    and under $\Q^S:=\Q^N$ the log-price has drift $r-y+\frac12|\sigma_S|^2$ instead of
    $\nu=r-y-\frac12|\sigma_S|^2$.
  \item (\emph{Currency put--call symmetry.}) In a currency market as in \cref{sec:ma-fx},
    with no assumption on the dynamics of $X$, let $N_t=X_te^{r_ft}$ (the foreign money market
    in domestic currency) be a positive price process and $\Q_f=\Q^N$, with $r=r_d$. A payoff of
    $H\ge0$ units of foreign currency at $T$ is worth $e^{-r_fT}\E^{\Q_f}[H]$ in foreign currency
    at $0$, and for $K>0$
    \[
      e^{-r_dT}\EQ\big[(X_T-K)^+\big]=X_0K\cdot e^{-r_fT}\E^{\Q_f}\Big[\Big(\frac1K-\frac1{X_T}\Big)^+\Big]:
    \]
    a call on one unit of foreign currency struck at $K$ is worth $X_0K$ times the foreign
    price of a put on one unit of domestic currency struck at $1/K$. In the model
    \eqref{eq:ma-gk}, $1/X$ has drift $r_f-r_d$ and volatility $\sigma$ under $\Q_f$, so the
    put is the Garman--Kohlhagen put with the roles of $r_d$ and $r_f$ exchanged.
\end{enumerate}
\end{proposition}

\begin{proof}
(i) $Z_T>0$, and $\EQ Z_T=1$ because $N/B$ is a martingale; equivalence follows from
$Z_T>0$. The martingale property gives $\EQ[Z_T\mid\F_t]=N_tB_0/(N_0B_t)=:Z_t$.

(ii) By Bayes' rule for conditional expectations (\textcite{shreve2}, Lemma~5.2.2), for $s\le T$
and $\F_s$-measurable $Y$ with $\E^{\Q^N}|Y|<\infty$,
$\E^{\Q^N}[Y\mid\F_t]=\EQ[Z_sY\mid\F_t]/Z_t$ for $t\le s$. With $Y=V_s/N_s$, $Z_sY=(B_0/N_0)V_s/B_s$, so
\[
  \E^{\Q^N}\Big[\frac{V_s}{N_s}\,\Big|\,\F_t\Big]=\frac{B_t}{N_t}\,\EQ\Big[\frac{V_s}{B_s}\,\Big|\,\F_t\Big],
  \qquad \E^{\Q^N}\Big|\frac{V_s}{N_s}\Big|=\frac{B_0}{N_0}\,\EQ\Big|\frac{V_s}{B_s}\Big| .
\]
Hence $V/N$ is $\Q^N$-integrable if and only if $V/B$ is $\Q$-integrable, and
$V_t/N_t=\E^{\Q^N}[V_s/N_s\mid\F_t]$ if and only if $V_t/B_t=\EQ[V_s/B_s\mid\F_t]$. The pricing
identity is the display with $s=T$ and $V_T=H$.

(iii) The density process $Z_t=\exp(\theta\cdot W_t-\frac12|\theta|^2t)$ is the one of
Girsanov's theorem (\textcite{karatzas-shreve}, Theorem~3.5.1), so $W^N$ is a $\Q^N$-Brownian
motion, and substituting $\dd W_t=\dd W^N_t+\theta\,\dd t$ gives the dynamics. The covariation of
$\log S$ and $\log N$ is $\sigma_S\cdot\theta\,\dd t$, since $\dd\log N_t=(r-\frac12|\theta|^2)\,\dd t+\theta\cdot\dd W_t$.

(iv) $N_t/B_t=S_te^{(y-r)t}=S_0\exp(\sigma_S\cdot W_t-\frac12|\sigma_S|^2t)$, so $\theta=\sigma_S$
and $Z_T=S_Te^{-(r-y)T}/S_0=S_T/\EQ[S_T]$. By (iii) the drift of $\dd S/S$ becomes
$r-y+|\sigma_S|^2$, and that of $\log S$ is $r-y+\frac12|\sigma_S|^2$.

(v) A foreign payoff $H$ is the domestic payoff $X_TH$; by (ii), which holds for nonnegative payoffs
whether or not their prices are finite (Bayes' rule does), its domestic price is
$N_0\E^{\Q_f}[X_TH/N_T]=X_0e^{-r_fT}\E^{\Q_f}[H]$, which is $e^{-r_fT}\E^{\Q_f}[H]$ units of foreign
currency. The identity follows from $(X_T-K)^+=X_TK(1/K-1/X_T)^+$, a foreign payoff of
$K(1/K-1/X_T)^+$. In the model \eqref{eq:ma-gk}, $N_t/B_t=X_0\exp(\sigma W_t-\frac12\sigma^2t)$, so by
(iii) $\dd X/X=(r_d-r_f+\sigma^2)\,\dd t+\sigma\,\dd W^N$ under $\Q_f$, and It\^o's formula gives
$\dd(1/X)/(1/X)=(r_f-r_d)\,\dd t+\sigma\,\dd(-W^N)$, with $-W^N$ a $\Q_f$-Brownian motion.
\end{proof}

The hypothesis that $N/B$ is a true martingale matters: a positive local martingale is a
supermartingale, and if $\EQ[N_T/B_T]<N_0/B_0$ then $\Q^N$ has total mass below one. In (iii)
and (iv) the exponential has constant coefficients, and it is a martingale.
\textcite{geman-elkaroui-rochet} developed the method in this generality; \textcite{shreve2},
Section~9.2 (Theorem~9.2.2) and Sections~9.3.3 and 9.3.7, gives it for the foreign measure and the
put--call duality of (v). Under the share measure $N(d_+)$ is the probability that the call
finishes in the money, as noted after \cref{thm:bsf-formula}, and the density in (iv) is the one used by
\cref{rem:is-share}.

\begin{example}[The course's London trader]\label{exa:ma-london}
A trader in London counts in pounds. For her the dollar is the foreign currency, worth
$X'_0=1/1.2=0.833333$ GBP, with rate $r_{\mathrm{GBP}}=2.5\%$ and yield $r_{\mathrm{USD}}=3\%$, and by (v) the cable
call of \cref{exa:ma-cable} is $K=1.206015$ dollar puts struck at $1/K$ pounds. The
Garman--Kohlhagen put is worth $0.0354710$ GBP, so the call is worth $K\times0.0354710=0.0427786$
GBP, the course's value, which is $0.0513343$ USD at the spot rate. Her hedge is the same
portfolio: the put's delta, $-0.463940$ dollars per put, times $K$ gives $-0.559519$ USD, and the
rest is $0.509044$ GBP in cash, as in \cref{exa:ma-cable}. Both traders hold the same assets,
pounds and dollars; only the unit of account differs.
\end{example}

The two forms of put--call symmetry are one identity. In the model \eqref{eq:ma-gk}, $X_T/F$
has under $\Q_f$ the law that $F/X_T$ has under $\Q$ (both log-normal with log-mean
$\frac12\sigma^2T$ and log-variance $\sigma^2T$), and substituting this into (v) turns the foreign put
into $K/F$ domestic puts struck at $F^2/K$: this is \cref{cor:sr-pcs}, whose applications to static
hedging of barriers are discussed there.

\begin{remark}[Monte Carlo in another num\'eraire]\label{rem:ma-mc-numeraire}
By (ii), $V_0=N_0\,\E^{\Q^N}[H/N_T]$: one may simulate under $\Q^N$ and average $N_0H/N_T$. This
is importance sampling with the likelihood ratio $\dd\Q/\dd\Q^N=N_0B_T/(N_TB_0)$
(\cref{prop:is-likelihood-ratio}). Under the share measure the call becomes
$S_0e^{-yT}\,\E^{\Q^S}\big[(1-K/S_T)^+\big]$, an expectation of a payoff between $0$ and $1$. In
the test case, with $2^{20}$ draws (our implementation, \texttt{numpy} \texttt{default\_rng},
seed $3$), the standard deviation per draw is $9.95$ under $\Q^S$ against $12.57$ under $\Q$, a
variance ratio of $1.6$; \exref{ex:ma-share} shows that the first stays bounded as $\sigma$ grows while the
second does not.
\end{remark}

\section{Exchange options}\label{sec:ma-exchange}

From here on, $d$ assets follow, under $\Q$,
\begin{equation}\label{eq:ma-model}
  \frac{\dd S_{i,t}}{S_{i,t}}=(r-y_i)\,\dd t+\sigma_i\,\dd W_{i,t},\qquad
  \dd\langle W_i,W_j\rangle_t=\rho_{ij}\,\dd t,\qquad i,j=1,\dots,d,
\end{equation}
with $\sigma_i>0$ and a correlation matrix $\rho=(\rho_{ij})$, symmetric positive semidefinite with
unit diagonal. Concretely, $W=A\bar W$ for a standard $k$-dimensional Brownian motion $\bar W$ and
a $d\times k$ matrix $A$ with $AA^\top=\rho$ (\cref{prop:rng-cholesky}), and $(\F_t)$ is generated by $\bar W$. The log-prices
$\log S_{i,T}=\log S_{i,0}+\nu_iT+\sigma_iW_{i,T}$, $\nu_i=r-y_i-\frac12\sigma_i^2$, form a Gaussian
vector with $\E S_{i,T}=F_i:=S_{i,0}e^{(r-y_i)T}$ and covariance
$\Sigma_{ij}=\rho_{ij}\sigma_i\sigma_jT$.

\begin{theorem}[Margrabe's formula]\label{thm:ma-margrabe}
Let $d=2$, $\rho=\rho_{12}\in[-1,1]$, $T>0$, and
\[
  \sigma_{\mathrm{ex}}=\sqrt{\sigma_1^2-2\rho\sigma_1\sigma_2+\sigma_2^2}\ \in\ \big[|\sigma_1-\sigma_2|,\,\sigma_1+\sigma_2\big].
\]
\begin{enumerate}[label=(\roman*)]
  \item The option to exchange asset $2$ for asset $1$ at $T$, paying $(S_{1,T}-S_{2,T})^+$, is
    worth, at $0$,
    \[
      M=\Black\big(S_{1,0}e^{-y_1T},\,S_{2,0}e^{-y_2T},\,\sigma_{\mathrm{ex}}^2T\big)
      =\mathrm{BS}\big(S_{1,0},S_{2,0},T;\,y_2,\,y_1,\,\sigma_{\mathrm{ex}}\big):
    \]
    the Black--Scholes call with spot $S_{1,0}$, strike $S_{2,0}$, rate $y_2$, yield $y_1$ and
    volatility $\sigma_{\mathrm{ex}}$. The interest rate $r$ does not enter.
  \item $\sigma_{\mathrm{ex}}=0$ if and only if $\rho=1$ and $\sigma_1=\sigma_2$; then
    $M=(S_{1,0}e^{-y_1T}-S_{2,0}e^{-y_2T})^+$. In all cases $\sigma_{\mathrm{ex}}$ decreases
    strictly in $\rho$, from $\sigma_1+\sigma_2$ at $\rho=-1$ to $|\sigma_1-\sigma_2|$ at $\rho=1$.
  \item If $\sigma_{\mathrm{ex}}>0$, $M$ is homogeneous of degree one in $(S_{1,0},S_{2,0})$ with
    $\partial M/\partial S_{1,0}=e^{-y_1T}N(d_+)$ and $\partial M/\partial S_{2,0}=-e^{-y_2T}N(d_-)$,
    $d_\pm=\big(\log\frac{S_{1,0}e^{-y_1T}}{S_{2,0}e^{-y_2T}}\pm\frac12\sigma_{\mathrm{ex}}^2T\big)/(\sigma_{\mathrm{ex}}\sqrt T)$,
    and $M=S_{1,0}\,\partial_1M+S_{2,0}\,\partial_2M$: the delta hedge holds the two assets and no cash.
\end{enumerate}
\end{theorem}

\begin{proof}
(ii) $\sigma_{\mathrm{ex}}^2=(\sigma_1-\sigma_2)^2+2(1-\rho)\sigma_1\sigma_2$, which vanishes exactly
when both terms do, and is affine in $\rho$ with slope $-2\sigma_1\sigma_2<0$.

(i) Write $W_i=a_i\cdot\bar W$, with $a_i$ the $i$-th row of $A$, so that $|a_i|=1$ and
$a_1\cdot a_2=\rho$; this holds for every $k$, including the rank-one factors at $\rho=\pm1$. The
num\'eraire $N_t=e^{y_2t}S_{2,t}$ has $N_t/B_t=S_{2,0}\exp(\sigma_2a_2\cdot\bar W_t-\frac12\sigma_2^2t)$, so it
satisfies \cref{prop:ma-numeraire}(iii), with $\bar W$ in the role of $W$ and $\theta=\sigma_2a_2$. Under
$\Q^N$, $\bar W^N=\bar W-\theta t$ is a standard Brownian motion; the drift of $\dd S_1/S_1$ rises by
$\sigma_1a_1\cdot\theta=\rho\sigma_1\sigma_2$ and that of $\dd S_2/S_2$ by $\sigma_2^2$. Hence
$R=S_1/S_2$ satisfies
\[
  \log R_T=\log R_0+\big(y_2-y_1-\tfrac12\sigma_{\mathrm{ex}}^2\big)T+u\cdot\bar W^N_T,\qquad
  u=\sigma_1a_1-\sigma_2a_2,
\]
because $\rho\sigma_1\sigma_2-\frac12\sigma_1^2-\frac12\sigma_2^2=-\frac12\sigma_{\mathrm{ex}}^2$; and
$|u|^2=\sigma_1^2-2\rho\sigma_1\sigma_2+\sigma_2^2=\sigma_{\mathrm{ex}}^2$. So under $\Q^N$, $R_T$ is lognormal with mean $R_0e^{(y_2-y_1)T}$ and
log-variance $\sigma_{\mathrm{ex}}^2T$ (a constant if $\sigma_{\mathrm{ex}}=0$). By
\cref{prop:ma-numeraire}(ii) with $H=(S_{1,T}-S_{2,T})^+$, for which $\EQ H\le\EQ S_{1,T}<\infty$,
\begin{align*}
  M&=N_0\,\E^{\Q^N}\Big[\frac{(S_{1,T}-S_{2,T})^+}{e^{y_2T}S_{2,T}}\Big]
  =S_{2,0}e^{-y_2T}\,\E^{\Q^N}\big[(R_T-1)^+\big]\\
  &=S_{2,0}e^{-y_2T}\Black\big(R_0e^{(y_2-y_1)T},1,\sigma_{\mathrm{ex}}^2T\big),
\end{align*}
and homogeneity of $\Black$ gives the first form; the second is the identity after
\eqref{eq:ma-black} with $r=y_2$, $y=y_1$. For $\sigma_{\mathrm{ex}}=0$ this is the formula in (ii).

(iii) The derivatives are $\Delta$ and $\partial_Kc$ of \cref{prop:bsf-greeks} for the inputs of (i), and
the last identity is Euler's theorem (\cref{prop:sr-euler}).
\end{proof}

With $y_1=y_2=0$ this is the formula of \textcite{margrabe}. In the course's words, the
num\'eraire change is a \emph{dimensional reduction}: writing the price as $S_2\,D(S_1/S_2)$ turns
the two-asset pricing equation into a one-asset equation for $D$, the price measured in units of
asset $2$. The same device is the pre-wash of \cref{prop:sr-power}. Options on the better or the
worse of two assets reduce to exchange options (\exref{ex:ma-max}); the spread
$(S_{1,T}-S_{2,T}-K)^+$ with $K\ne0$ does not (\exref{ex:ma-spread}).

\paragraph{Cross rates and the volatility triangle.} For a dollar investor, the pound and the euro
are two assets of the form \eqref{eq:ma-model}: $S_1=X_1$, the USD price of a pound, with yield
$r_{\mathrm{GBP}}$, and $S_2=X_2$, the USD price of a euro, with yield $r_{\mathrm{EUR}}$. The cross rate $X_{12}=X_1/X_2$
is the EUR price of a pound, and a GBP/EUR call struck at $K$ pays $(X_{12,T}-K)^+$ euros, that is
$(X_{1,T}-KX_{2,T})^+$ dollars: an exchange option. Its volatility is
$\sigma_{12}=\sigma_{\mathrm{ex}}=\big(\sigma_1^2-2\rho\sigma_1\sigma_2+\sigma_2^2\big)^{1/2}$, with $\rho$ the
correlation of GBP/USD and EUR/USD returns, both quoted against the dollar. The three
volatilities are the sides of a triangle with angle $\arccos\rho$ opposite $\sigma_{12}$ (the law of
cosines), so $|\sigma_1-\sigma_2|\le\sigma_{12}\le\sigma_1+\sigma_2$, and a quoted cross volatility
implies the correlation $\rho=(\sigma_1^2+\sigma_2^2-\sigma_{12}^2)/(2\sigma_1\sigma_2)$. Dividing the dollar
price by $X_{2,0}$ gives the Garman--Kohlhagen price of a euro-based investor with volatility
$\sigma_{12}$: the law of one price at work.

\begin{example}[The course's cross-currency inputs]\label{exa:ma-cross}
With $X_1=1.20$, $X_2=1.02$, $\sigma_1=11\%$, $\sigma_2=10\%$, $r_{\mathrm{USD}}=3\%$, $r_{\mathrm{GBP}}=2.5\%$ and
$r_{\mathrm{EUR}}=1\%$, the GBP/EUR spot is $1.176471$ and its one-year forward is
$1.176471\,e^{0.01-0.025}=1.158955$. Correlations in $[-1,1]$ allow
$\sigma_{12}\in[1\%,21\%]$. The ATMF GBP/EUR call, paying $(1\text{ GBP}-1.158955\text{ EUR})^+$, is
worth between $0.004669$ and $0.097871$ USD; a quoted $\sigma_{12}=8\%$ implies $\rho=0.7136$ and
the value $0.037343$ USD (\cref{fig:ma-cross}). A quote outside $[1\%,21\%]$ admits no
correlation and is an arbitrage within the model; \exref{ex:ma-triangle} asks what a static
portfolio of four options (the cross call, the GBP/USD call, and the EUR/USD call and put) can do with it.
\end{example}

\begin{figure}[ht]
  \centering
  \includegraphics[width=\textwidth]{ch17-cross}
  \caption{The cross rate GBP/EUR with $\sigma_1=11\%$ (GBP/USD), $\sigma_2=10\%$ (EUR/USD) and the
    rates of \cref{exa:ma-cross}. (a) The cross volatility against the correlation $\rho$ of
    GBP/USD and EUR/USD, between the dotted levels $|\sigma_1-\sigma_2|$ and $\sigma_1+\sigma_2$; the dot is
    the quote $8\%$, implying $\rho=0.714$. (b) The USD value of the
    ATMF GBP/EUR call by \cref{thm:ma-margrabe}, with Monte Carlo estimates and $\pm2$ standard
    errors, from the same $2^{14}$ normal pairs at every $\rho$ (seed $17$) and the principal-components
    factor, which also works at $\rho=\pm1$.}
  \label{fig:ma-cross}
\end{figure}

\section{Correlated simulation and quantos}\label{sec:ma-quanto}

\paragraph{Simulating correlated assets.} Under \eqref{eq:ma-model}, a vector of terminal prices is
sampled exactly in one step,
\[
  S_{i,T}=S_{i,0}\exp\big(\nu_iT+\sigma_i\sqrt T\,(AZ)_i\big),\qquad Z\sim\mathcal N(0,I_k),\quad AA^\top=\rho,
\]
by \cref{prop:rng-cholesky}(i), and a path is sampled by repeating this over each interval with
independent $Z$'s (\cref{prop:ps-gbm-exact}). The Cholesky factor costs about $d^3/3$ operations
once and $d^2/2$ per draw; a lower triangular factor with positive diagonal exists if and only if
$\rho$ is positive definite. At
$\rho_{12}=\pm1$ the second pivot $1-\rho_{12}^2$ vanishes and the positive-pivot recursion stops, while the
principal-components factor of \cref{prop:rng-cholesky}(iii), which needs only semidefiniteness,
still works; we use it in \cref{fig:ma-cross}. A matrix of pairwise correlations collected from
different sources need not be semidefinite at all, and then no model has it. Antithetic draws
negate the whole vector $Z$; moment matching must match the sample covariance as well as the
means, which is harder (\cref{prop:cv-moment-matching}).

\paragraph{Quantos.} A \emph{quanto} pays a foreign quantity in domestic currency at a rate fixed in
advance: an option on a foreign index, struck in foreign currency, whose payoff is converted at the
rate $\bar X$, whatever the exchange rate at maturity. Let $S$ be the foreign-currency price of a
foreign stock with dividend yield $y$, $X$ the exchange rate as in \cref{sec:ma-fx}, and $\Q$ the
domestic risk-neutral measure: the domestic money market, the foreign money market $X_te^{r_ft}$
and the foreign stock with reinvested dividends, worth $X_tS_te^{yt}$, all have $\Q$-martingale
prices once discounted at $r_d$. Assume
\begin{equation}\label{eq:ma-quanto-model}
  \frac{\dd X_t}{X_t}=\mu_X\,\dd t+\sigma_X\,\dd W_{X,t},\qquad
  \frac{\dd S_t}{S_t}=\mu_S\,\dd t+\sigma_S\,\dd W_{S,t},\qquad
  \dd\langle W_X,W_S\rangle_t=\rho\,\dd t,
\end{equation}
with constants $\mu_X,\mu_S$, $\sigma_X,\sigma_S>0$ and $\rho\in[-1,1]$, where $W_X=a_X\cdot\bar W$ and
$W_S=a_S\cdot\bar W$ for a standard $k$-dimensional $\Q$-Brownian motion $\bar W$ generating the filtration
and unit vectors with $a_X\cdot a_S=\rho$, as in \eqref{eq:ma-model}. As in \cref{def:sm-model}, $\Q$ is
equivalent to the real-world measure $\Prob$ on $\F_T$.

\begin{proposition}[Quantos and composites]\label{prop:ma-quanto}
In this market, with $T>0$ and $K>0$:
\begin{enumerate}[label=(\roman*)]
  \item $\mu_X=r_d-r_f$ and $\mu_S=r_f-y-\rho\sigma_S\sigma_X$. Under the foreign measure $\Q_f$ of
    \cref{prop:ma-numeraire}(v), $\dd S/S=(r_f-y)\,\dd t+\sigma_S\,\dd W^f_S$ with $W^f_S$ a
    $\Q_f$-Brownian motion.
  \item (\emph{Quanto.}) The payoff $\bar X(S_T-K)^+$ in domestic currency costs
    \[
      \bar Xe^{-r_dT}\Black\big(F^q,K,\sigma_S^2T\big)=\bar X\,\mathrm{BS}(S_0,K,T;\,r_d,\,y_q,\,\sigma_S),\qquad
      F^q=S_0e^{(r_f-y-\rho\sigma_S\sigma_X)T},
    \]
    with $y_q=y+r_d-r_f+\rho\sigma_S\sigma_X$.
  \item (\emph{Composite.}) The payoff $(X_TS_T-K)^+$, a foreign stock struck in domestic currency,
    costs $\mathrm{BS}(X_0S_0,K,T;\,r_d,\,y,\,\sigma_c)$ with
    $\sigma_c^2=\sigma_S^2+2\rho\sigma_S\sigma_X+\sigma_X^2$.
  \item (\emph{Foreign strike, spot conversion.}) The payoff $X_T(S_T-K)^+$ costs
    $X_0\,\mathrm{BS}(S_0,K,T;\,r_f,\,y,\,\sigma_S)$.
\end{enumerate}
\end{proposition}

\begin{proof}
(i) By It\^o's formula, $X_te^{(r_f-r_d)t}$ and $X_tS_te^{(y-r_d)t}$ are exponentials of
Brownian motions with constant coefficients, with drift rates $\mu_X+r_f-r_d$ and
$\mu_X+\mu_S+\rho\sigma_S\sigma_X+y-r_d$ (the cross term is $\dd\langle\log X,\log S\rangle=\rho\sigma_S\sigma_X\,\dd t$).
Such an exponential is a martingale if and only if its drift rate is zero (its mean is
$e^{(\text{drift rate})\,t}$ times its initial value), which gives $\mu_X$ and then $\mu_S$. Under $\Q_f$,
\cref{prop:ma-numeraire}(iii) with $N_t/B_t=X_0\exp(\sigma_XW_{X,t}-\frac12\sigma_X^2t)$, that is
$\theta=\sigma_Xa_X$, raises the drift of $\dd S/S$ by $\sigma_Sa_S\cdot\theta=\rho\sigma_S\sigma_X$.
(ii) Under $\Q$, $S_T$ is lognormal with mean $F^q$ and log-variance $\sigma_S^2T$; and
$\mathrm{BS}(S_0,K,T;r_d,y_q,\sigma_S)=e^{-r_dT}\Black(S_0e^{(r_d-y_q)T},K,\sigma_S^2T)$ with $r_d-y_q=r_f-y-\rho\sigma_S\sigma_X$.
(iii) By (i), $\log(X_TS_T)$ is normal with variance $\sigma_c^2T$, and $\EQ[X_TS_T]=X_0S_0e^{(r_d-y)T}$
because $X_tS_te^{(y-r_d)t}$ is a martingale. (iv) By \cref{prop:ma-numeraire}(v) the foreign-currency
price is $e^{-r_fT}\E^{\Q_f}[(S_T-K)^+]=\mathrm{BS}(S_0,K,T;r_f,y,\sigma_S)$ by (i); convert at $X_0$.
\end{proof}

The \emph{quanto adjustment} $-\rho\sigma_S\sigma_X$ lowers the drift of a foreign asset that tends to
rise with the foreign currency. The correlation is that of the returns of $S$, in foreign
currency, and of $X$, quoted as domestic per foreign. Quadratic covariation is the limit in
probability of sums of products of increments, and limits in probability are the same under
equivalent measures; since $\Prob$, $\Q$ and $\Q_f$ are equivalent, $\rho$ is the same under all three, and it
can be estimated from historical data; quoting the rate the other way round changes its sign. The
adjustment does not vanish when $r_d=r_f$, and it is the only trace of the exchange rate in the
price: the quanto does not depend on $X_0$, and $\sigma_X$ enters only through $\rho\sigma_X$.
\exref{ex:ma-elf} prices the remaining contract of the course's list, a currency option whose
notional is a foreign share.

\begin{example}[A quanto call]\label{exa:ma-quanto}
Take a foreign index with $S_0=K=100$ GBP, $y=2\%$ and $\sigma_S=20\%$, the cable parameters of
\cref{exa:ma-cable}, and $\rho=-0.3$. Then $y_q=1.84\%$, the quanto forward is $101.1668$ against
the foreign forward $100.5013$, and the quanto call with $\bar X=1$ USD per GBP costs $8.35441$
USD; with $\rho=0$ and $\rho=+0.3$ it would cost $7.99516$ and $7.64664$. The ordinary call costs
$8.03523$ GBP, that is $9.64228$ USD at the spot rate (iv), and the composite call struck at $120$
USD costs $9.79110$ USD, with $\sigma_c=19.72\%$ (iii). Our implementation, simulating $(S_T,X_T)$
under $\Q$ with $2^{20}$ draws (\texttt{numpy} \texttt{default\_rng}, seed $2026$), gives $8.3457$, $9.6313$
and $9.7784$, with standard errors between $0.013$ and $0.016$, all within one standard error of the formulas.
\end{example}

\section{Basket options}\label{sec:ma-basket}

A \emph{basket call} pays $(A-K)^+$, where $A=\sum_iw_iS_{i,T}$ is a weighted sum of the prices
at maturity in the model \eqref{eq:ma-model}: an option on an index of a few stocks, or on a
currency basket. It is the cross-sectional analogue of the Asian option of
\cref{ch:15-asian}: there the summands are one asset at several dates, with covariances
$\sigma^2\min(t_i,t_j)\ge0$; here they are several assets at one date, with the covariances
$\Sigma_{ij}=\rho_{ij}\sigma_i\sigma_jT$ of any sign. The sum of lognormals is not lognormal, and
with $d\ge2$ assets no closed form is known, except in degenerate cases. Write
$D=e^{-rT}$, $Y=D(A-K)^+$, $C=\EQ Y$ and $\Phi=\sum_iw_iF_i=\EQ A$.

\begin{proposition}[Basket options]\label{prop:ma-basket}
\begin{enumerate}[label=(\roman*)]
  \item (\emph{The geometric basket.}) Let all $w_i>0$, $a_i=w_iF_i/\Phi$, and
    $G=\Phi\prod_i(S_{i,T}/F_i)^{a_i}$. Then $\log G\sim\mathcal N(\mu_G,v_G)$ with
    $\mu_G=\log\Phi-\frac12\sum_ia_i\Sigma_{ii}$ and $v_G=a^\top\Sigma a$; $G\le A$ on every path;
    and $X^{\mathrm g}=D(G-K)^+$ has mean $D\Black(\EQ G,K,v_G)\le C$, where
    $\EQ G=e^{\mu_G+v_G/2}\le\Phi$.
  \item (\emph{The geometric control.}) Under (i), $0\le Y-X^{\mathrm g}\le D(A-G)$ and, with
    $c=\Sigma a$,
    \begin{align*}
      \Var(Y-X^{\mathrm g})&\le D^2\EQ\big[(A-G)^2\big]\\
      &=D^2\Big(\sum_{i,j}w_iw_jF_iF_je^{\Sigma_{ij}}-2\,\EQ G\sum_iw_iF_ie^{c_i}+(\EQ G)^2e^{v_G}\Big).
    \end{align*}
    If $v_G>0$ and $\Var Y>0$, the control-variate estimator with the optimal coefficient has variance
    $(1-\Corr(Y,X^{\mathrm g})^2)\Var Y/n$ (\cref{thm:cv-optimal}), and
    $1-\Corr(Y,X^{\mathrm g})^2\le D^2\EQ[(A-G)^2]/\Var Y$, whose numerator does not depend on $K$.
  \item (\emph{Moment matching.}) $\EQ A=\Phi$ and $\EQ A^2=\sum_{i,j}w_iw_jF_iF_je^{\Sigma_{ij}}$. If
    $\Phi>0$, the lognormal law with these two moments has log-variance
    $v_2=\log(\EQ A^2/\Phi^2)$, and the \emph{two-moment approximation} is
    $C_2=D\Black(\Phi,K,v_2)$. It equals $C$ if $A$ is lognormal (for instance if $d=1$, or if
    $\rho_{ij}=1$ and all $\sigma_i$ are equal), and if all $w_i\ge0$ and $K\le0$. It is not a bound:
    in the three-asset example below, $C_2-C$ changes sign with $K$ (\cref{tab:ma-basket}).
  \item (\emph{Mixed signs.}) If some weights are positive and some negative and $\Sigma$ is
    positive definite, then $\Q(A<0)>0$. The geometric basket of (i)--(ii) is then undefined, and no
    lognormal law has the law of $A$: when $\Phi>0$, $C_2$ can still be computed, but it puts no mass where
    $A<0$, and none of the exact cases of (iii) applies; if $\Phi\le0$, no lognormal law even has the
    mean of $A$. Plain Monte Carlo is unaffected.
\end{enumerate}
\end{proposition}

\begin{proof}
(i) $\log G=\log\Phi+\sum_ia_i(\log S_{i,T}-\log F_i)$ is a linear function of a Gaussian vector,
with $\E(\log S_{i,T}-\log F_i)=-\frac12\Sigma_{ii}$ and variance $a^\top\Sigma a$. The weights $a_i$ are
positive with sum one, and $\log$ is concave, so the weighted arithmetic--geometric mean
inequality gives $\prod_ix_i^{a_i}\le\sum_ia_ix_i$ for $x_i=S_{i,T}/F_i$; multiplied by $\Phi$, this
is $G\le\sum_iw_iS_{i,T}=A$. The mean of $X^{\mathrm g}$ is the identity after \eqref{eq:ma-black}
(or $D(\EQ G-K)^+$ if $v_G=0$); it is at most $C$ because $(G-K)^+\le(A-K)^+$, and
$\EQ G\le\EQ A$.
(ii) $x\mapsto(x-K)^+$ is nondecreasing and $1$-Lipschitz, so $0\le(A-K)^+-(G-K)^+\le A-G$, and
$\Var(Y-X^{\mathrm g})\le\EQ[(Y-X^{\mathrm g})^2]$. For the moments, $\EQ[S_{i,T}S_{j,T}]=F_iF_je^{\Sigma_{ij}}$,
$\EQ G^2=(\EQ G)^2e^{v_G}$, and $\log S_{i,T}+\log G$ is normal with variance
$\Sigma_{ii}+2c_i+v_G$, where $c_i=\Cov(\log S_{i,T},\log G)=(\Sigma a)_i$, so
$\EQ[S_{i,T}G]=F_i\,\EQ G\,e^{c_i}$. If $v_G>0$ then $X^{\mathrm g}$ has positive variance, and
\cref{thm:cv-optimal}(ii) applies; its optimal variance is at most the variance at coefficient one.
(iii) The moments follow as in (ii). A lognormal $L$ with mean $\Phi$ has $\E L^2=\Phi^2e^{\Var\log L}$,
which gives $v_2$. If $A$ is lognormal, $C_2=C$ by \eqref{eq:ma-black}; under the stated
condition all $S_{i,T}/F_i$ are equal to one lognormal variable, so $A=\Phi\,S_{1,T}/F_1$. If all
$w_i\ge0$ and $K\le0$, then $A-K\ge0$ and both sides equal $D(\Phi-K)$. The sign changes are in
\cref{tab:ma-basket}.
(iv) The vector $(\log S_{i,T})$ has a positive Gaussian density on $\R^d$, so every open subset of
$(0,\infty)^d$ has positive probability. With $w_1>0>w_2$, say, the open set where
$w_1S_{1,T}+\sum_{i\ge3}|w_i|S_{i,T}<|w_2|S_{2,T}$ is nonempty and lies in $\{A<0\}$. A lognormal
variable is positive, and $G>0$.
\end{proof}

The geometric basket with the weights $a_i$, proportional to the forward values $w_iF_i$ rather
than to $w_i$, is the choice of \textcite{gentle-basket}, as it is usually described: $A/\Phi$ and $G/\Phi$ are then the
arithmetic and geometric means of the same ratios $S_{i,T}/F_i$, which agree to first order when the
ratios are close to one. Equal weights $w_i/\sum_jw_j$ would give a geometric mean of the prices
themselves, which is far below $A$ when the prices differ in scale. Gentle's \emph{approximation}
prices $(G-K')^+$ with the strike $K'=K-(\Phi-\EQ G)$ shifted by the difference of the means; like
$C_2$, which is due to \textcite{levy-1992}, it is a heuristic, with no error bound, and it errs in both
directions: it underprices in \cref{exa:ma-basket}, but with three equicorrelated assets ($\rho_{ij}=0.9$,
volatilities $5\%$, $5\%$, $80\%$, $T=2$, $K=60$, other inputs as in \cref{exa:ma-basket}) it gives $38.265$
against $C=37.343$, and $C_2=40.992$. The bound in (ii)
does not depend on the strike, as in \cref{prop:as-control}, but it depends on how dispersed the
ratios $S_{i,T}/F_i$ are, and dispersion grows with low or negative correlations
(\exref{ex:ma-dispersion}).

\begin{remark}[Conditioning, and a reference price]\label{rem:ma-reference}
Write the log-prices as $m+LZ$ with $LL^\top=\Sigma$ and $Z\sim\mathcal N(0,I_d)$, and split
$Z=e\,\xi+QZ'$, where $e$ is the unit vector along $L^\top a$ and the columns of $Q$ are an orthonormal
basis of $e^\perp$, so that $\log G-\mu_G=\sqrt{v_G}\,\xi$, with
$\xi\sim\mathcal N(0,1)$ independent of $Z'\in\R^{d-1}$. Given $Z'$, $A=\sum_i\alpha_ie^{b_i\xi}$ with
$\alpha_i>0$ and $b=Le=c/\sqrt{v_G}$; it is a convex function of $\xi$, so it exceeds $K$ outside an
interval, whose endpoints are found by bisection, and $\EQ[(A-K)^+\mid Z']$ is a sum of terms
$e^{b_i^2/2}N(\cdot)$, as for the Asian options of \cref{ch:15-asian}. This conditional
value is an unbiased estimator with no larger variance (\cref{thm:st-conditional}) and a smooth function
of $Z'$ when all $b_i>0$. In the example below ($d=3$), a product Gauss--Hermite rule in $Z'$ with
$8\times8$ nodes already gives ten digits, and $16\times16$ nodes agree with $64\times64$ to $10^{-14}$.
Conditioning on $G$ instead gives the lower bound $D\,\EQ[(\EQ[A\mid G]-K)^+]\le C$ by the conditional
Jensen inequality, whatever the covariances, with
$\EQ[A\mid\log G=x]=\sum_iw_iF_i\exp\big(c_i(x-\mu_G)/v_G-c_i^2/(2v_G)\big)$. If every $c_i\ge0$, this function is
increasing in $x$, it crosses $K$ once, and the bound has the closed form of \cref{prop:as-curran}(ii), as
\textcite{curran} proposed for ``portfolio options''.
\end{remark}

\begin{example}[A three-asset basket]\label{exa:ma-basket}
Take three assets with the test-case inputs $S_{i,0}=100$, $r=1\%$, $y_i=2\%$, $T=1$, volatilities
$20\%$, $25\%$ and $30\%$, the correlations $0.5$, $0.3$, $0.6$ of \cref{exm:rng-three-assets}, equal
weights $w_i=\frac13$ and $K=100$. Then $\Phi=99.0050$, $\sqrt{v_G}=20.23\%$ and $\sqrt{v_2}=20.30\%$,
against volatilities averaging $25\%$: diversification. The reference price of
\cref{rem:ma-reference} is $C=7.4655807$. Our implementation, with $2^{20}$ draws (\texttt{numpy}
\texttt{default\_rng}, seed $2026$), gives
$7.46603\pm0.01256$ by plain Monte Carlo; the geometric control, with the in-sample coefficient
$1.0498$ and correlation $0.99679$, gives $7.46541\pm0.00101$, $12$ times more precise; and three
controls (the geometric call, the forward $DA$ and the strip $D\sum_iw_i(S_{i,T}-K)^+$) give
$7.46585\pm0.00071$, with $R^2=0.99678$ (\cref{thm:cv-multiple}). Before any simulation, (ii) bounds
$\Var(Y-X^{\mathrm g})$ by $2.888$, about $1.7\%$ of $\Var Y\approx165.5$. The two-moment
approximation gives $C_2=7.48004$, Gentle's $7.36489$, the conditioning bound of
\cref{rem:ma-reference} $7.45406$, and the geometric call, a lower bound, $6.87976$.
\end{example}

\begin{table}[ht]
  \centering
  \small
  \begin{tabular}{@{}lrrrrr@{}}
    \toprule
    $K$ & $80$ & $90$ & $100$ & $110$ & $120$\\
    \midrule
    reference $C$ & $20.132481$ & $12.804691$ & $7.465581$ & $4.023289$ & $2.027888$\\
    two moments $C_2$ & $+319.2$ & $+319.7$ & $+144.6$ & $-68.7$ & $-199.0$\\
    Gentle & $-325.6$ & $-700.3$ & $-1006.9$ & $-1085.2$ & $-945.9$\\
    conditioning bound & $-76.8$ & $-105.0$ & $-115.2$ & $-110.5$ & $-95.5$\\
    geometric call & $-9647.6$ & $-7881.5$ & $-5858.2$ & $-3991.3$ & $-2517.6$\\
    \bottomrule
  \end{tabular}
  \caption{The three-asset basket call of \cref{exa:ma-basket}: reference prices from
    \cref{rem:ma-reference}, and the errors of the approximations and bounds, in units of $10^{-4}$.}
  \label{tab:ma-basket}
\end{table}

\begin{figure}[ht]
  \centering
  \includegraphics[width=\textwidth]{ch17-basket}
  \caption{The three-asset basket call of \cref{exa:ma-basket} against the strike. (a) Standard
    deviation per draw: plain Monte Carlo, the geometric control and the three controls, all with
    in-sample coefficients, from the same $2^{16}$ draws (seed $17$). (b) Errors relative to the reference
    price of the two-moment approximation, of Gentle's approximation and of the conditioning lower bound;
    the geometric call, $0.25$ to $0.96$ below the reference, is off the scale.}
  \label{fig:ma-basket}
\end{figure}

\Cref{tab:ma-basket} and \cref{fig:ma-basket} show the pattern of \cref{ch:15-asian} with larger
numbers. The control removes between $98.8\%$ and $99.6\%$ of the variance for strikes from $80$ to
$120$, less as the option moves out of the money. Two-moment matching errs by up to $0.03$ with
both signs, overpricing in the money and underpricing out of it; Gentle's shift underprices by
$0.03$ to $0.11$. For a basket the summands are less alike than the fixings of an average, so
approximations that treat $A$ as nearly lognormal, or nearly geometric, are less accurate, and the
control is weaker: $1-\rho^2$ is $6.4\times10^{-3}$ here against $6.9\times10^{-4}$ for the monthly
Asian of \cref{ch:15-asian}. Beyond three or four assets, product quadrature in $Z'$ becomes too
expensive (\cref{prop:mc-dimension}), while Monte Carlo with the geometric control, or the conditional
estimator of \cref{rem:ma-reference} with randomised quasi-Monte Carlo, keeps its cost per draw of
order $d^2$.

\demo{17-multi-asset.ipynb}{the cable call in either currency, with simulation under the domestic and
the foreign measures; the cross-rate triangle and exchange options against the correlation;
quantos against the correlation; and the basket's control variates and approximations, with sliders
for the volatilities, the correlations and the strike.}

\begin{incode}
\coderef{code/src/mfe230d/multiasset.py}: \texttt{black} is \eqref{eq:ma-black};
\texttt{garman\_kohlhagen}, \texttt{gk\_hedge} and \texttt{currency\_symmetry} implement
\cref{prop:ma-garman-kohlhagen} and \cref{prop:ma-numeraire}(v), on top of
\coderef{code/src/mfe230d/bs.py}; \texttt{numeraire\_drift} and \texttt{call\_samples\_share\_measure}
\cref{prop:ma-numeraire}(iii)--(iv) and \cref{rem:ma-mc-numeraire}; \texttt{terminal\_prices} samples
\eqref{eq:ma-model} with \texttt{rng.cholesky} or \texttt{rng.pca\_factor}; \texttt{margrabe},
\texttt{exchange\_vol}, \texttt{implied\_correlation} and \texttt{cross\_atmf\_call}
\cref{thm:ma-margrabe} and \cref{exa:ma-cross}; \texttt{quanto\_option}, \texttt{composite\_option},
\texttt{foreign\_option} and \texttt{fx\_equity\_terminal} \cref{prop:ma-quanto};
\texttt{geometric\_basket\_option}, \texttt{basket\_cv\_bound}, \texttt{basket\_lognormal},
\texttt{gentle\_call} and \texttt{basket\_mc} (through \texttt{variance.control\_variate} and
\texttt{variance.regression\_control\_variates}) \cref{prop:ma-basket}; \texttt{basket\_reference} and
\texttt{curran\_lower\_bound} \cref{rem:ma-reference}.
\end{incode}

\paragraph{Course examples reproduced.} The ATMF cable call ($0.05133$ USD, $0.04278$ GBP) and its
two deltas ($0.5090$ GBP, $-0.5595$ USD); the London trader's thought experiment, with the same
price and the same hedge in pounds, and put--call symmetry; the cross-rate volatility
$\sigma_{\mathrm{eff}}^2=\sigma_1^2-2\rho\sigma_1\sigma_2+\sigma_2^2$, the Frankfurt trader's
thought experiment and the triangle of admissible volatilities; the Cholesky recipe for correlated
normals (via \cref{exm:rng-three-assets}); and the course's ``pre-washed'' inputs for quantos,
composites and foreign options.

\furtherreading
\textcite{shreve2}, Chapter~9, treats num\'eraires, the domestic and foreign risk-neutral measures,
Siegel's paradox (Section~9.3.4) and put--call duality, and \textcite{geman-elkaroui-rochet} is the
original account of the change of num\'eraire. \textcite{garman-kohlhagen} and \textcite{margrabe} are
short and readable; \textcite{stulz-1982} prices options on the maximum or the minimum of two assets
with the bivariate normal distribution. \textcite{glasserman-mc} samples multivariate normals in
Section~2.3.3 and correlated geometric Brownian motions in Section~3.2.3, and treats control variates in
Section~4.1. \textcite{levy-1992}, \textcite{gentle-basket} and \textcite{curran} give the basket
approximations of \cref{sec:ma-basket}.
